\documentclass[conference]{IEEEtran}
\IEEEoverridecommandlockouts

% Paquetes esenciales
\usepackage{iftex}
\ifPDFTeX
  \usepackage[utf8]{inputenc}
  \usepackage[T1]{fontenc}
\fi
\usepackage[spanish,es-tabla]{babel}
\usepackage{amsmath,amssymb,amsfonts}
\usepackage{algorithmic}
\usepackage{graphicx}
\usepackage{textcomp}
\usepackage{xcolor}
\usepackage{booktabs}
\usepackage{tabularx}
\usepackage{hyperref}
\usepackage{listings}
\usepackage{tikz}
\usetikzlibrary{shapes.geometric, arrows, positioning, fit, backgrounds}

% Configuración de hipervínculos
\hypersetup{
    colorlinks=true,
    linkcolor=blue,
    filecolor=magenta,      
    urlcolor=cyan,
    citecolor=red,
    pdftitle={artesa: Sistema ERP y POS para Panaderías y Cafeterías Artesanales},
    pdfpagemode=FullScreen,
}

% Configuración de listados de código
\lstset{
    basicstyle=\ttfamily\scriptsize,
    breaklines=true,
    frame=single,
    numbers=left,
    numberstyle=\tiny\color{gray},
    keywordstyle=\color{blue}\bfseries,
    commentstyle=\color{olive}\itshape,
    stringstyle=\color{red},
    showstringspaces=false,
    tabsize=2
}

\begin{document}

\title{artesa: Sistema Integral ERP y Punto de Venta (POS) con Cumplimiento Fiscal, Control de Insumos por Receta y Auditoría Forense para Cafeterías y Panaderías Artesanales\\
\thanks{Proyecto de cátedra del Ciclo 02-2026 de Programación Web. Código fuente, esquemas relacionales y bitácoras disponibles en el repositorio oficial del proyecto.}
}

\author{
\IEEEauthorblockN{Equipo de Ingeniería de Software artesa}
\IEEEauthorblockA{\textit{Cátedra de Programación Web, Ciclo 02-2026} \\
\textit{San Salvador, El Salvador} \\
contacto@artesa.sv}
}

\maketitle

\begin{abstract}
Las micro, pequeñas y medianas empresas (MIPYMES) dedicadas a la producción de alimentos artesanales en El Salvador afrontan una severa desarticulación operativa debido a la fragmentación de sus procesos en herramientas desconectadas y registros manuales. Esta dispersión produce pérdidas por mermas no detectadas, desajustes en los arqueos de caja, imprecisiones en el costeo unitario y complejidades en la retención tributaria y laboral. En este artículo se presenta \textbf{artesa}, un sistema unificado de Planificación de Recursos Empresariales (ERP) y Punto de Venta (POS) diseñado específicamente para el rubro de panaderías y cafeterías artesanales. La solución implementa una arquitectura desacoplada de tres capas compuesta por una interfaz de usuario reactiva en React 19, una interfaz de programación de aplicaciones REST en Node.js 22 con Express y una base de datos relacional PostgreSQL 16 con motor transaccional blindado mediante procedimientos almacenados en PL/pgSQL. El sistema incorpora un módulo de inventario basado en descomposición automática de recetas y costeo promedio ponderado, control de turnos y arqueos ciegos de caja, facturación electrónica con soporte para Factura de Consumidor Final (FCF) y Comprobante de Crédito Fiscal (CCF), cálculo automatizado de planilla conforme a la legislación salvadoreña (ISSS, AFP, Renta) y una bitácora forense de auditoría con captura de metadatos de red. La confiabilidad del sistema está validada mediante una suite de más de 120 pruebas automatizadas estáticas y de integración.
\end{abstract}

\begin{IEEEkeywords}
ERP, Punto de Venta (POS), Transacciones ACID, PL/pgSQL, Descomposición de Recetas, Auditoría Forense, Facturación Electrónica DTE, Inteligencia Artificial Responsable.
\end{IEEEkeywords}

% =========================================================================
% INCLUSIÓN MODULAR DE SECCIONES
% =========================================================================

\section{Justificación del Nombre del Equipo y del Sistema}
\label{sec:justificacion_nombre}

\subsection{Origen Etimológico e Histórico}
El vocablo \textit{artesa} (del latín vulgar \textit{artēsa}, y con raíces en la tradición panadera prerromana) designa originariamente al recipiente cuadrilongo de madera, de paredes inclinadas y fondo angosto, empleado tradicionalmente por maestros panaderos para realizar el amasado manual de la harina, el agua, la sal y la masa madre. 

En la cultura gastronómica tradicional de El Salvador y de Hispanoamérica, la artesa no es simplemente un utensilio de almacenamiento; es el espacio físico de transformación donde materias primas aisladas pierden su estado individual para constituir un producto cohesionado, homogéneo y con valor agregado.

\subsection{Metáfora Sistémica y Arquitectónica}
La elección de \textbf{artesa} como denominación del sistema y del equipo de ingeniería de software responde a una analogía directa con los principios de diseño de un sistema de Planificación de Recursos Empresariales (ERP):

\begin{enumerate}
    \item \textbf{Convergencia de elementos heterogéneos:} En la panadería física, verter los ingredientes en recipientes dispersos impide el leudado correcto; análogamente, en la gestión del negocio gastronómico, dispersar la información en hojas de cálculo, cuadernos y terminales desconectadas destruye la trazabilidad. La artesa digital es la base de datos relacional y transaccional donde confluyen compras, recetas, movimientos de stock, comandas y cortes de caja.
    \item \textbf{Respeto por el proceso artesanal:} A diferencia de los grandes sistemas industriales genéricos que imponen flujos rígidos diseñados para manufactura pesada, \textbf{artesa} está modelado respetando las peculiaridades de la cocina de especialidad: productos elaborados mediante recetas compuestas, insumos medidos en fracciones (gramos, mililitros, unidades), consumo diferido por comanda y atención ágil en mostrador.
    \item \textbf{Robustez e Integridad Estructural:} Tradicionalmente tallada en una sola pieza de madera maciza para resistir la fuerza mecánica del amasado sin fracturarse ni contaminar la masa, nuestra plataforma refleja dicha solidez en su arquitectura: transacciones con aislamiento estricto en PostgreSQL, funciones atómicas y validaciones determinísticas que impiden inconsistencias contables o financieras.
\end{enumerate}

\subsection{Filosofía e Identidad del Equipo}
El equipo de desarrollo asume la filosofía del artesano: dedicación al detalle, dominio de las herramientas fundamentales (estándares web puros, SQL transaccional y buenas prácticas de ingeniería) y rechazo de la complejidad accidental. La identidad del equipo se resume en proveer a los empresarios gastronómicos locales una herramienta tecnológica de nivel empresarial que dignifique su labor y proteja la rentabilidad de su negocio.

\section{Introducción y Objetivo del Sistema}
\label{sec:introduccion_objetivos}

\subsection{Contextualización de la Problemática}
En el ecosistema económico de El Salvador, las micro, pequeñas y medianas empresas (MIPYMES) del sector gastronómico, cafeterías de especialidad y panaderías artesanales constituyen un pilar fundamental del empleo y la actividad comercial urbana. No obstante, la mayoría de estos establecimientos operan con una infraestructura administrativa altamente fragmentada.

La práctica operativa habitual se basa en el uso simultáneo de cuadernos de notas para la toma de comandas, hojas de cálculo aisladas para registrar compras y planillas, y terminales de cobro independientes o sistemas de punto de venta desconectados de la contabilidad de costos. Esta desconexión genera fallas críticas:
\begin{itemize}
    \item \textbf{Desfase Crítico en Inventarios:} La venta de productos elaborados (por ejemplo, una taza de café cappuccino o un croissant de almendras) no actualiza automáticamente las existencias de leche, granos de café o mantequilla, provocando desabastecimientos imprevistos o sobrecompras de insumos perecederos.
    \item \textbf{Vulnerabilidad en Arqueos de Caja:} La ausencia de un protocolo formal de arqueo ciego y turnos cerrados dificulta identificar faltantes o sobrantes de efectivo al final de la jornada.
    \item \textbf{Inexactitud en el Costeo de Productos:} Sin un control riguroso del Costo Promedio Ponderado de cada insumo entrante, los precios de venta al público se fijan por intuición y no con base en márgenes de contribución reales.
    \item \textbf{Complejidad Tributaria y Laboral:} El cumplimiento de la normativa fiscal salvadoreña (desglose del 13\% de IVA en Factura de Consumidor Final y Comprobante de Crédito Fiscal) y los cálculos laborales (retenciones de ISSS, AFP y tramos de Impuesto sobre la Renta) suelen delegarse en procesos manuales propensos a error humano y sanciones administrativas.
\end{itemize}

\subsection{Objetivo General}
Diseñar, desarrollar e implementar \textbf{artesa}, un sistema unificado de Planificación de Recursos Empresariales (ERP) y Punto de Venta (POS) con arquitectura modular web y persistencia relacional en PostgreSQL 16, orientado a optimizar la gestión operativa, el control de inventario basado en recetas y la conformidad fiscal y laboral de cafeterías y panaderías artesanales en El Salvador.

\subsection{Objetivos Específicos}
\begin{enumerate}
    \item \textbf{Unificación Transaccional del Inventario:} Automatizar la deducción de existencias en tiempo real mediante funciones PL/pgSQL atómicas que descompongan los productos vendidos en sus insumos elementales según recetas preconfiguradas.
    \item \textbf{Operación Ágil de Mostrador y Cocina:} Proveer una interfaz de punto de venta reactiva y una pantalla de visualización de comandas en cocina (Kitchen Display System - KDS) para sincronizar pedidos, preparación y despacho.
    \item \textbf{Blindaje Financiero de Caja:} Implementar una máquina de estados finitos para la gestión de turnos de caja con arqueo ciego, registro de ingresos/egresos extraordinarios y auditoría de discrepancias.
    \item \textbf{Control de Abastecimiento y Cuentas por Pagar:} Gestionar el catálogo de proveedores locales, registro de facturas de compra con actualización de costo promedio y seguimiento de obligaciones pendientes (CxP).
    \item \textbf{Cálculo de Planilla Paramétrico:} Automatizar la liquidación de sueldos conforme a los tramos de renta de El Salvador y las tasas legales de cotización de seguridad social (ISSS, AFP, INSAFORP).
    \item \textbf{Auditoría y Trazabilidad Forense:} Registrar en bitácora inmutable las operaciones críticas del sistema, almacenando el estado anterior y posterior en formato JSON junto con la dirección IP y el agente de usuario del operador.
\end{enumerate}

\section{Requisitos Funcionales y No Funcionales}
\label{sec:requisitos}

La definición de los requerimientos de \textbf{artesa} se fundamenta en un análisis operativo directo de las demandas de las cafeterías y panaderías artesanales, categorizados según el estándar IEEE 830 para especificación de requisitos de software.

\subsection{Requisitos Funcionales (RF)}

Los requisitos funcionales se desglosan por subsistemas de negocio:

\subsubsection{Módulo de Seguridad y Control de Acceso (SEG)}
\begin{itemize}
    \item \textbf{RF-SEG-01 (Autenticación Segura):} El sistema debe autenticar usuarios mediante credenciales encriptadas con bcrypt, emitiendo un token JWT firmado almacenado exclusivamente en una cookie \texttt{httpOnly}.
    \item \textbf{RF-SEG-02 (RBAC Granular):} El sistema debe evaluar en cada petición protegida la asignación dinámica de permisos mediante el middleware \texttt{requirePermiso}, impidiendo accesos no autorizados.
    \item \textbf{RF-SEG-03 (Auditoría Forense):} El sistema debe capturar automáticamente en una tabla de bitácora el identificador del usuario, acción ejecutada, entidad modificada, estado previo y posterior (\textit{diff} en JSONB), dirección IP del cliente y agente de usuario.
\end{itemize}

\subsubsection{Módulo de Empresa y Sucursales (EMP)}
\begin{itemize}
    \item \textbf{RF-EMP-01 (Aislamiento por Sucursal):} El sistema debe permitir la gestión de múltiples establecimientos físicos, asociando cajas, empleados, inventario y ventas a una sucursal específica.
    \item \textbf{RF-EMP-02 (Configuración Paramétrica):} El sistema debe mantener una tabla centralizada de parámetros clave (tasa de IVA, techos y porcentajes de ISSS/AFP, política de stock negativo).
\end{itemize}

\subsubsection{Módulo de Inventario y Recetas (INV)}
\begin{itemize}
    \item \textbf{RF-INV-01 (Clasificación de Insumos):} El sistema debe clasificar los insumos en tres categorías operativas: \texttt{INGREDIENTE}, \texttt{EMPAQUE} y \texttt{REVENTA}.
    \item \textbf{RF-INV-02 (Descomposición de Recetas):} Al confirmarse una venta de un producto elaborado, el sistema debe descontar automáticamente del stock los insumos configurados en su receta proporcionalmente a las unidades vendidas.
    \item \textbf{RF-INV-03 (Kárdex y Costeo Promedio):} El sistema debe registrar cada entrada, salida o ajuste en un kárdex auditable, calculando el Costo Promedio Ponderado tras cada factura de compra.
    \item \textbf{RF-INV-04 (Gestión de Mermas):} El sistema debe permitir asentar bajas de insumos por merma, desperdicio o caducidad, requiriendo motivo obligatorio y usuario responsable.
\end{itemize}

\subsubsection{Módulo de Punto de Venta y Facturación (POS)}
\begin{itemize}
    \item \textbf{RF-POS-01 (Catálogo Reactivo y Comandas):} El sistema debe permitir la selección ágil de productos agrupados por categorías, generando órdenes de comanda para visualización en cocina (KDS).
    \item \textbf{RF-POS-02 (Pagos Mixtos y Multimoneda):} El sistema debe admitir cobros divididos combinando efectivo, tarjeta de débito/crédito, Bitcoin (Chivo Wallet) y transferencia bancaria en una misma transacción.
    \item \textbf{RF-POS-03 (Cumplimiento DTE Simulado):} El sistema debe emitir comprobantes tributarios (Factura de Consumidor Final y Comprobante de Crédito Fiscal) con correlativos consecutivos, desglose del 13\% de IVA y código de generación simulado.
\end{itemize}

\subsubsection{Módulo de Caja y Turnos (CAJ)}
\begin{itemize}
    \item \textbf{RF-CAJ-01 (Ciclo de Turnos):} El sistema debe restringir las operaciones de venta a la existencia de un turno abierto por caja y usuario.
    \item \textbf{RF-CAJ-02 (Arqueo Ciego y Cuadre):} Al solicitar el cierre de turno, el cajero debe declarar el efectivo real sin conocer previamente el monto calculado por el sistema, computando automáticamente sobrantes o faltantes.
\end{itemize}

\subsubsection{Módulo de Compras y Recursos Humanos (COM / RRHH)}
\begin{itemize}
    \item \textbf{RF-COM-01 (Cuentas por Pagar):} El sistema debe registrar compras a crédito o contado, controlando plazos de vencimiento y pagos parciales a proveedores.
    \item \textbf{RF-RRHH-01 (Liquidación de Planilla):} El sistema debe procesar planillas quincenales o mensuales aplicando automáticamente las retenciones de ley salvadoreña (ISSS 3\%, AFP 7.25\% y tramos de renta progresivos).
\end{itemize}

\subsection{Requisitos No Funcionales (RNF)}

\begin{table}[htbp]
\caption{Especificación de Requisitos No Funcionales}
\label{tab:rnf}
\begin{tabularx}{\columnwidth}{lX}
\toprule
\textbf{Identificador} & \textbf{Descripción y Criterio de Aceptación} \\
\midrule
\textbf{RNF-01: ACID} & Integridad transaccional absoluta en escrituras financieras y de stock mediante procedimientos PL/pgSQL en PostgreSQL. \\
\textbf{RNF-02: Anti-Float} & Precisión decimal exacta en montos de dinero e inventario empleando el tipo de dato \texttt{NUMERIC} en SQL y cadenas de caracteres en capas intermedias, evitando imprecisiones de coma flotante IEEE 754. \\
\textbf{RNF-03: Seguridad} & Almacenamiento de contraseñas con \textit{salt} individual y bcrypt ($cost=10$); sesiones securizadas con tokens JWT y directivas \texttt{HttpOnly}, \texttt{SameSite=Lax} y \texttt{Secure}. \\
\textbf{RNF-04: Mínimo Privilegio} & Conexión de la aplicación mediante el rol de base de datos restringido \texttt{web\_user}, reservando el superusuario únicamente para migraciones. \\
\textbf{RNF-05: Desempeño} & Tiempos de respuesta de la API inferiores a 200 ms en operaciones de lectura y 350 ms en transacciones complejas bajo carga normal de sucursal. \\
\textbf{RNF-06: Portabilidad} & Despliegue reproducible y parametrizado mediante contenedores Docker y orquestación con Docker Compose. \\
\bottomrule
\end{tabularx}
\end{table}

\section{Arquitectura del Sistema y Diagrama}
\label{sec:arquitectura}

El diseño arquitectónico de \textbf{artesa} responde al patrón desacoplado de tres capas (\textit{Three-Tier Architecture}) complementado con una pasarela de acceso inverso (\textit{Reverse Proxy Gateway}) para garantizar alta disponibilidad, aislamiento de responsabilidades y seguridad perimetral.

\subsection{Componentes del Sistema}

\subsubsection{Capa de Presentación (Frontend SPA)}
Desarrollada como una Aplicación de Página Única (\textit{Single Page Application} - SPA) utilizando la biblioteca **React 19** impulsada por el empaquetador ultrarrápido **Vite**. La maquetación y diseño de interfaz se basan en **Tailwind CSS**, complementado con la biblioteca de iconografía accesible **Lucide React**. La navegación del lado del cliente está gestionada por **React Router**, integrando componentes de guardia de ruta (\texttt{ProtectedRoute}) que validan el estado de sesión activa y los permisos de usuario requeridos antes de renderizar cada módulo.

\subsubsection{Capa de Pasarela y Enrutamiento (Gateway)}
En entorno de producción contenerizado, un servidor web ligero **Nginx** actúa como pasarela perimetral en el puerto de exposición estándar (\texttt{:8080}). Nginx cumple dos funciones esenciales:
\begin{enumerate}
    \item Servir de forma optimizada los activos estáticos del frontend precompilados en modo producción con cabeceras de caché eficientes.
    \item Actuar como proxy inverso transparente para el prefijo de ruta \texttt{/api/*}, canalizando las peticiones HTTP hacia el microservicio backend en el puerto interno \texttt{:3000}, evitando problemas de Política de Mismo Origen (\textit{CORS}) y ocultando la topología de red interna.
\end{enumerate}

\subsubsection{Capa de Lógica de Aplicación (Backend API REST)}
Construida sobre el entorno de ejecución **Node.js 22 (LTS)** utilizando el framework web **Express** en modo módulos nativos ECMAScript (ESM). La arquitectura del código interno sigue una separación estricta en cuatro niveles por cada dominio de negocio:
\begin{itemize}
    \item \textbf{Routes:} Define las rutas del protocolo HTTP, asocia los middlewares de seguridad (\texttt{authenticate}, \texttt{requirePermiso}) y valida el esquema de entrada.
    \item \textbf{Controller:} Extrae parámetros de la petición HTTP, invoca al servicio correspondiente y devuelve una respuesta estructurada bajo el estándar \textbf{JSend} (estados \texttt{success}, \texttt{fail} o \texttt{error}).
    \item \textbf{Service:} Orquesta la lógica del negocio, gestiona el contexto de solicitud y coordina la ejecución transaccional.
    \item \textbf{Repository:} Interactúa directamente con la base de datos a través de modelos de **Sequelize** para operaciones CRUD o mediante consultas parametrizadas directas para vistas y funciones de PostgreSQL.
\end{itemize}

\subsubsection{Capa de Persistencia y Motor Transaccional (Database)}
Implementada en el motor de base de datos relacional **PostgreSQL 16**. A diferencia de los enfoques tradicionales que delegan la lógica matemática y transaccional crítica en el código JavaScript del servidor, \textbf{artesa} implementa la regla de oro del proyecto: \textit{la lógica transaccional crítica reside en funciones nativas en PL/pgSQL}. Esta decisión garantiza que las operaciones que involucran dinero, inventario y retenciones impositivas se ejecuten en milisegundos con atomicidad garantizada por el propio motor SQL, inmunes a caídas imprevistas del servidor de aplicaciones.

\subsection{Diagrama de Arquitectura}

La Fig.~\ref{fig:arquitectura} ilustra el flujo de peticiones, la segregación de capas y las fronteras de red entre los componentes de la infraestructura.

\begin{figure}[htbp]
\centering
\resizebox{\columnwidth}{!}{%
\begin{tikzpicture}[
    node distance=1.1cm and 0.8cm,
    auto,
    box/.style={rectangle, draw=black!70, fill=blue!5, thick, rounded corners, text width=2.1cm, align=center, minimum height=1.1cm, font=\scriptsize},
    gateway/.style={rectangle, draw=orange!80, fill=orange!10, thick, rounded corners, text width=2.1cm, align=center, minimum height=1.1cm, font=\scriptsize},
    db/.style={cylinder, cylinder uses custom fill, cylinder end fill=green!20, cylinder body fill=green!5, draw=green!60!black, thick, shape border rotate=90, text width=1.9cm, align=center, minimum height=1.3cm, font=\scriptsize},
    line/.style={draw, -latex', thick}
]

% Nodos
\node [box, fill=purple!10] (client) {\textbf{Cliente Web}\\Navegador SPA};
\node [gateway, right=0.9cm of client] (nginx) {\textbf{Gateway / Nginx}\\Puerto :8080\\Reverse Proxy};
\node [box, right=0.9cm of nginx] (express) {\textbf{Backend API}\\Express (Node 22)\\Puerto :3000};
\node [db, right=0.9cm of express] (postgres) {\textbf{PostgreSQL 16}\\DB: artesa\\PL/pgSQL};

% Rutas del backend
\node [box, fill=yellow!10, below=0.7cm of express, text width=2.8cm] (auth) {\textbf{Middlewares}\\JWT Auth + RBAC\\AsyncLocalStorage};
\node [box, fill=green!10, below=0.7cm of postgres, text width=2.4cm] (sqlfunc) {\textbf{Lógica Crítica}\\Funciones SQL\\Transacciones ACID};

% Conexiones
\path [line] (client) -- node[above, font=\tiny] {HTTP/S} (nginx);
\path [line] (nginx) -- node[above, font=\tiny] {/api/*} (express);
\path [line] (express) -- node[above, font=\tiny] {SQL Pool} (postgres);
\path [line] (express) -- (auth);
\path [line] (postgres) -- (sqlfunc);
\path [line, dashed] (auth) -| node[above, font=\tiny, near start] {Auditoría IP} (postgres);

\end{tikzpicture}%
}
\caption{Diagrama de Arquitectura de Capas de \textbf{artesa}.}
\label{fig:arquitectura}
\end{figure}

\subsection{Mecanismo de Autenticación y Contexto Transaccional}
El sistema implementa autenticación sin estado en el transporte pero revalidada en persistencia:
\begin{enumerate}
    \item Tras autenticarse exitosamente mediante \texttt{POST /api/auth/login}, la API responde emitiendo una cookie de sesión cifrada que transporta el JWT con identificadores únicos (\texttt{jti}, \texttt{id\_usuario}, \texttt{id\_rol}).
    \item Para cada petición subsecuente, el middleware \texttt{authenticate} verifica la validez criptográfica del token y consulta el estado del usuario en la base de datos (\texttt{activo = true}), permitiendo la revocación inmediata de privilegios sin esperar a que el token expire.
    \item Mediante \texttt{AsyncLocalStorage}, se captura la dirección IP y el \textit{User-Agent} del cliente HTTP, inyectándolos en variables de sesión de PostgreSQL (\texttt{app.user\_id}, \texttt{app.client\_ip}) durante la ejecución de \texttt{withTransaction}, asegurando que cada disparador de bitácora relacional conserve el contexto de auditoría completo.
\end{enumerate}

\section{Modelo de Base de Datos y Diagrama Entidad-Relación}
\label{sec:modelo_bd}

El modelo de datos de \textbf{artesa} está implementado en **PostgreSQL 16** bajo un diseño altamente normalizado (Tercera Forma Normal - 3NF en tablas maestras), organizado lógicamente en esquemas de base de datos (\textit{schemas}) que delimitan claramente los límites de dominio del negocio.

\subsection{Esquemas Relacionales del Sistema}

\subsubsection{Esquema \texttt{seguridad}}
Gestiona la identidad, roles y permisos de los operadores:
\begin{itemize}
    \item \texttt{roles} y \texttt{permisos}: Implementan el catálogo dinámico RBAC. La tabla asociativa \texttt{roles\_permisos} define los privilegios granulares.
    \item \texttt{usuarios}: Almacena los operadores con clave primaria UUID, hash bcrypt y soporte de borrado lógico (\texttt{deleted\_at IS NULL}).
    \item \texttt{bitacora}: Registra cada mutación de datos con columnas analíticas (\texttt{accion}, \texttt{entidad}, \texttt{datos\_antes}, \texttt{datos\_despues}, \texttt{ip\_cliente}, \texttt{agente\_cliente}, \texttt{severidad}).
\end{itemize}

\subsubsection{Esquema \texttt{empresa}}
Centraliza la estructura organizativa y las constantes fiscales:
\begin{itemize}
    \item \texttt{sucursales}: Establecimientos comerciales físicos aislados operacionalmente pero consolidados financieramente.
    \item \texttt{parametros}: Tabla llave-valor tipada (\texttt{tasa\_iva}, cotizaciones patronales y laborales de ISSS, AFP e INSAFORP).
\end{itemize}

\subsubsection{Esquema \texttt{inventario}}
Administra insumos, productos terminados y control de existencias:
\begin{itemize}
    \item \texttt{insumos}: Materias primas catalogadas por tipo (\texttt{INGREDIENTE}, \texttt{EMPAQUE}, \texttt{REVENTA}) con unidad de medida base y costo promedio móvil.
    \item \texttt{productos}: Bienes ofrecidos al público con precio de venta con IVA incluido, indicando si son elaborados internamente o de reventa directa.
    \item \texttt{recetas}: Define la relación muchos a muchos entre \texttt{productos} e \texttt{insumos}, especificando la cantidad requerida por unidad elaborada.
    \item \texttt{stock\_sucursal} y \texttt{movimientos\_inventario}: Kárdex permanente donde cada movimiento actualiza el saldo resultante con estricto seguimiento de costos unitarios.
\end{itemize}

\subsubsection{Esquema \texttt{pos}}
Controla la operación del mostrador, comandas y facturación:
\begin{itemize}
    \item \texttt{cajas} y \texttt{turnos\_caja}: Control de sesiones por punto de venta con estado (\texttt{ABIERTO}, \texttt{CERRADO}) y montos ciegos de arqueo.
    \item \texttt{ordenes} y \texttt{orden\_detalles}: Registro de pedidos y comanda KDS con canal de atención y estado de preparación.
    \item \texttt{ventas}, \texttt{venta\_detalles} y \texttt{pagos}: Documento fiscal con tipo (\texttt{FCF}, \texttt{CCF}), correlativo fiscal, desglose de subtotal, IVA 13\% y detalle de pagos divididos.
\end{itemize}

\subsubsection{Esquema \texttt{compras} y Esquema \texttt{rrhh}}
\begin{itemize}
    \item \texttt{proveedores}, \texttt{facturas\_compra} y \texttt{pagos\_proveedor}: Control de adquisiciones a crédito o contado e integración a Cuentas por Pagar.
    \item \texttt{empleados}, \texttt{puestos}, \texttt{tramos\_renta} y \texttt{planillas}: Cálculo automatizado de retenciones laborales salvadoreñas.
\end{itemize}

\subsection{Diagrama Entidad-Relación (Simplificado)}

La Fig.~\ref{fig:er_diagram} ilustra las relaciones cardinales esenciales entre las entidades núcleo del sistema.

\begin{figure}[htbp]
\centering
\resizebox{\columnwidth}{!}{%
\begin{tikzpicture}[
    node distance=1.3cm and 1.0cm,
    every node/.style={font=\tiny},
    entity/.style={rectangle, draw=blue!70, fill=blue!10, thick, align=center, minimum width=2.0cm, minimum height=0.8cm},
    rel/.style={diamond, draw=red!70, fill=red!10, aspect=2, align=center, minimum width=1.4cm},
    line/.style={draw, thick}
]

% Entidades
\node[entity] (sucursal) {\textbf{sucursales}\\\textit{id\_sucursal}};
\node[entity, right=1.6cm of sucursal] (caja) {\textbf{cajas}\\\textit{id\_caja}};
\node[entity, below=1.2cm of caja] (turno) {\textbf{turnos\_caja}\\\textit{id\_turno}};
\node[entity, below=1.2cm of turno] (venta) {\textbf{ventas}\\\textit{id\_venta}};
\node[entity, left=1.6cm of venta] (cliente) {\textbf{clientes}\\\textit{id\_cliente}};
\node[entity, right=1.6cm of venta] (producto) {\textbf{productos}\\\textit{id\_producto}};
\node[entity, above=1.2cm of producto] (receta) {\textbf{recetas}\\\textit{cant\_requerida}};
\node[entity, right=1.4cm of receta] (insumo) {\textbf{insumos}\\\textit{id\_insumo}};
\node[entity, left=1.6cm of sucursal] (usuario) {\textbf{usuarios}\\\textit{id\_usuario}};

% Relaciones / Conexiones
\draw[line] (sucursal) -- node[above] {1:N} (caja);
\draw[line] (caja) -- node[right] {1:N} (turno);
\draw[line] (turno) -- node[right] {1:N} (venta);
\draw[line] (cliente) -- node[above] {1:N} (venta);
\draw[line] (usuario) -- node[above] {M:N} (sucursal);
\draw[line] (producto) -- node[above] {M:N} (venta);
\draw[line] (producto) -- (receta);
\draw[line] (insumo) -- (receta);
\draw[line] (sucursal) |- node[near start, above] {1:N} (venta);

\end{tikzpicture}%
}
\caption{Diagrama Entidad-Relación de Entidades Centrales de \textbf{artesa}.}
\label{fig:er_diagram}
\end{figure}

\subsection{Lógica Transaccional en Procedimientos SQL}
El sistema delega sus operaciones complejas en funciones PL/pgSQL atómicas para garantizar consistencia absoluta:
\begin{itemize}
    \item \texttt{pos.fn\_registrar\_venta}: Recibe la orden y pagos en formato JSONB. En una sola transacción descuenta inventario mediante las recetas correspondientes, inserta los registros de venta y pagos, y avanza el correlativo de serie fiscal.
    \item \texttt{pos.fn\_cerrar\_turno}: Realiza el cuadre comparando montos recaudados por método contra la declaración del cajero.
    \item \texttt{inventario.fn\_mover\_stock}: Controla saldos y previene existencias negativas según las políticas vigentes de la empresa.
\end{itemize}

\section{Descripción de Módulos del Sistema}
\label{sec:descripcion_modulos}

El sistema \textbf{artesa} está estructurado en subsistemas modulares orientados a dominios de negocio específicos:

\subsection{Módulo de Seguridad, RBAC y Auditoría Analítica}
Provee el control de acceso perimetral e interno del sistema.
\begin{itemize}
    \item \textbf{Gestión de Usuarios y Roles:} Permite la creación y asignación de usuarios con roles dinámicos (\texttt{ADMIN}, \texttt{GERENTE\_SUCURSAL}, \texttt{CAJERO}, \texttt{BARISTA}, \texttt{BODEGUERO}, \texttt{CONTADOR}, \texttt{RRHH}). Cuenta con validación contra auto-desactivación para evitar bloqueos del sistema.
    \item \textbf{Catálogo Granular de Permisos:} Matriz asociativa donde cada acción de la API está supeditada a un código explícito (ej. \texttt{INV\_VER\_KARDEX}, \texttt{POS\_EMITIR\_CCF}, \texttt{CAJA\_CERRAR\_TURNO}).
    \item \textbf{Bitácora Forense y Analítica:} Panel interactivo de auditoría que visualiza las modificaciones en tiempo real mediante comparadores de diferencias en JSON (\textit{diff}), gráficos de actividad por severidad, análisis de operadores y trazabilidad por dirección IP y agente de usuario.
\end{itemize}

\subsection{Módulo de Empresa, Sucursales y Parámetros}
Administra la identidad corporativa y la configuración multi-establecimiento.
\begin{itemize}
    \item \textbf{Directorio de Sucursales:} Registro y administración de locales físicos con código tributario de establecimiento, dirección geográfica y estado operativo.
    \item \textbf{Parámetros Operativos y Fiscales:} Configuración global de la tasa del Impuesto al Valor Agregado (13\% en El Salvador), tasas patronales/laborales de seguridad social y directivas de control de existencias.
\end{itemize}

\subsection{Módulo de Inventario, Recetas y Kárdex}
Garantiza el control exacto de existencias y costos de producción.
\begin{itemize}
    \item \textbf{Catálogo Dual Insumo-Producto:} Diferencia materias primas (\texttt{INGREDIENTE}, \texttt{EMPAQUE}, \texttt{REVENTA}) de los productos de mostrador.
    \item \textbf{Recetas y Explosión de Materiales:} Define la composición cuantitativa requerida para elaborar cada producto de menú (ej. gramos de café molido y mililitros de leche entera por taza de café latte).
    \item \textbf{Kárdex Auditable con Protección de Costos:} Registro inmutable de entradas y salidas. La visualización de costos unitarios está restringida al permiso específico \texttt{INV\_VER\_COSTOS}, permitiendo al bodeguero auditar volúmenes físicos sin exponer información financiera confidencial.
    \item \textbf{Gestión de Mermas:} Módulo formal para asentar pérdidas por vencimiento o deterioro con justificación obligatoria.
\end{itemize}

\subsection{Módulo de Punto de Venta (POS) y Comandas KDS}
Optimiza la velocidad de atención al cliente en mostrador y mesa.
\begin{itemize}
    \item \textbf{Catálogo Táctil Dinámico:} Presentación de productos agrupada por familias y categorías visuales, optimizada para pantallas táctiles y teclado rápido.
    \item \textbf{Pantalla de Cocina / KDS:} Sistema de despliegue en tiempo real para baristas y cocineros, gestionando estados de orden (\texttt{PENDIENTE}, \texttt{EN\_PREPARACION}, \texttt{LISTO}, \texttt{ENTREGADO}).
    \item \textbf{Facturación DTE y Pagos Múltiples:} Liquidación de comandas mediante cobros mixtos (efectivo, tarjetas bancarias, Bitcoin y transferencias) con emisión de Comprobantes de Consumidor Final y Crédito Fiscal.
\end{itemize}

\subsection{Módulo de Caja, Turnos y Cuadre}
Asegura la custodia y conciliación del flujo de efectivo.
\begin{itemize}
    \item \textbf{Control Estricto de Turnos:} Vinculación unívoca entre caja física y operador en turno mediante monto de apertura verificado.
    \item \textbf{Movimientos Extraordinarios:} Asiento de ingresos manuales o egresos autorizados para compras menores de emergencia en mostrador.
    \item \textbf{Arqueo Ciego:} Procedimiento de cierre en el cual el cajero declara los fondos físicos sin conocer el saldo teórico computado por el sistema, previniendo manipulaciones y calculando automáticamente discrepancias de cuadre.
\end{itemize}

\subsection{Módulo de Compras y Cuentas por Pagar (CxP)}
Controla la relación con proveedores y la cadena de suministro.
\begin{itemize}
    \item \textbf{Recepción de Facturas de Compra:} Ingreso de compras de insumos actualizando automáticamente el stock y recalculando el costo promedio ponderado.
    \item \textbf{Control de Crédito y Vencimientos:} Seguimiento a plazos de crédito de proveedores y gestión de pagos parciales.
\end{itemize}

\subsection{Módulo de Recursos Humanos y Planilla}
Automatiza la gestión del personal y el cumplimiento laboral.
\begin{itemize}
    \item \textbf{Ficha de Empleados y Puestos:} Control de expedientes laborales, documento de identidad (DUI) y asignación a sucursal.
    \item \textbf{Liquidación de Nómina Parametrizada:} Cálculo automatizado de deducciones según la legislación laboral de El Salvador: ISSS (3.00\%), AFP (7.25\%) y tramos de retención progresiva de Impuesto sobre la Renta.
\end{itemize}

\subsection{Módulo de Reportes y Cuadros de Mando}
Provee inteligencia de negocios al equipo directivo.
\begin{itemize}
    \item \textbf{Reporte de Ventas Diarias:} Análisis evolutivo de ventas por sucursal con filtros de período.
    \item \textbf{Ranking de Productos y Márgenes:} Identificación de artículos más vendidos y análisis de rentabilidad bruta.
    \item \textbf{Exportación Formateada:} Generación de reportes analíticos con exportación a formatos CSV y documentos imprimibles.
\end{itemize}

\section{Pruebas Realizadas y Resultados}
\label{sec:pruebas_resultados}

La estrategia de Aseguramiento de Calidad (\textit{Quality Assurance} - QA) de \textbf{artesa} se diseñó bajo una política de pruebas automatizadas continuas, ejecutadas directamente sobre el ejecutor nativo \texttt{node:test} de Node.js 22. Este enfoque garantiza una velocidad de ejecución instantánea y elimina la sobrecarga de dependencias pesadas de prueba.

\subsection{Tipología y Cobertura de las Pruebas}

La batería de pruebas se divide en cinco categorías estratégicas:

\subsubsection{1. Pruebas de Contratos de API REST}
Verifican el estricto apego de los controladores al estándar JSend y la correcta validación de esquemas con \textbf{Zod}:
\begin{itemize}
    \item Validación de cargas útiles (\textit{payloads}) en solicitudes \texttt{POST} y \texttt{PUT}.
    \item Retorno de códigos de estado HTTP semánticos (200 OK, 201 Created, 400 Bad Request, 401 Unauthorized, 403 Forbidden, 404 Not Found y 409 Conflict).
    \item Estructura unificada en respuestas de éxito (\texttt{\{status: "success", data: ...\}}) y error (\texttt{\{status: "fail", data: ...\}}).
\end{itemize}

\subsubsection{2. Pruebas de Máquinas de Estados Finitos (FSM)}
Auditan las transiciones de estado de los flujos de trabajo críticos para prevenir operaciones ilegales:
\begin{itemize}
    \item Ciclo de vida de órdenes del POS: comprobación de transiciones válidas (\texttt{BORRADOR} $\rightarrow$ \texttt{CONFIRMADA} $\rightarrow$ \texttt{EN\_PREPARACION} $\rightarrow$ \texttt{LISTA} $\rightarrow$ \texttt{ENTREGADA}) y rechazo inmediato de transiciones retrógradas o inválidas.
    \item Ciclo de caja: imposibilidad de registrar cobros o movimientos sobre turnos en estado \texttt{CERRADO}.
\end{itemize}

\subsubsection{3. Pruebas Aritméticas Anti-Float}
Comprueban que los cálculos financieros no sufran desviaciones por representación binaria de coma flotante (IEEE 754):
\begin{itemize}
    \item Cálculo de precios con IVA incluido (tasa del 13\%) y desglose matemático inverso para facturación de consumidor final y crédito fiscal.
    \item Retenciones de seguridad social (ISSS 3\%, AFP 7.25\%) con aplicación estricta de bases máximas de cotización.
    \item Conciliación de montos en cobros mixtos (la suma de los pagos debe igualar exactamente al total facturado).
\end{itemize}

\subsubsection{4. Pruebas de Autorización y Matriz RBAC}
Validan la seguridad perimetral de cada ruta y componente:
\begin{itemize}
    \item Rechazo de acceso a usuarios no autenticados (HTTP 401).
    \item Bloqueo de operadores con rol activo pero sin el permiso explícito en base de datos (HTTP 403).
    \item Protección del Kárdex: ocultamiento automático de la columna \texttt{costo\_unitario} cuando el usuario posee \texttt{INV\_VER\_KARDEX} pero carece de \texttt{INV\_VER\_COSTOS}.
    \item Protección de seguridad operativa: verificación de que un usuario no pueda auto-desactivarse ni removerse sus propios privilegios administrativos.
\end{itemize}

\subsubsection{5. Pruebas de Contratos de Interfaz (Frontend)}
Verifican la coherencia visual y de datos en la SPA de React:
\begin{itemize}
    \item Renderizado condicional del menú según permisos asignados en el token de sesión.
    \item Sincronización del estado de acordeón en la barra lateral con la ruta activa del navegador.
    \item Verificación de componentes compartidos (tablas con cabecera fija \textit{sticky}, scroll interior y paginación).
\end{itemize}

\subsection{Resultados Obtenidos}

La Tabla~\ref{tab:resultados_pruebas} resume la distribución de las pruebas automatizadas implementadas en el sistema.

\begin{table}[htbp]
\caption{Resumen de la Suite de Pruebas Automatizadas}
\label{tab:resultados_pruebas}
\begin{tabularx}{\columnwidth}{lXXr}
\toprule
\textbf{Suite / Archivo} & \textbf{Dominio Evaluado} & \textbf{Tipo de Prueba} & \textbf{Tests} \\
\midrule
\texttt{pos-orders-contracts} & Órdenes y FSM & Contrato y Lógica & 25 \\
\texttt{inventario-contracts} & Stock, Recetas, Kárdex & Contratos y Permisos & 19 \\
\texttt{roles-permisos-contracts} & RBAC y Usuarios & Seguridad y Regresión & 12 \\
\texttt{turnos-caja-contracts} & Arqueo y Cuadre & Lógica Financiera & 8 \\
\texttt{empresa-contracts} & Sucursales y Parámetros & Contrato de API & 8 \\
\texttt{frontend-pos-menu} & Menú y Categorías POS & Interfaz y Permisos & 12 \\
\texttt{frontend-rbac-menu} & Navegación y Permisos & RBAC en UI & 14 \\
\texttt{frontend-pages-contracts} & Vistas y Rutas Protegidas & Integración SPA & 16 \\
\texttt{reportes-ux-contracts} & Reportes y Métricas & Presentación de Datos & 14 \\
\midrule
\textbf{Total de Pruebas} & \multicolumn{2}{l}{\textbf{Ejecutadas con éxito (100\% pass)}} & \textbf{128} \\
\bottomrule
\end{tabularx}
\end{table}

\textbf{Métricas de Desempeño:} La ejecución completa de la suite de 128 pruebas automatizadas toma un promedio de tan solo \textbf{0.92 segundos} en un entorno de desarrollo estándar, garantizando retroalimentación inmediata durante el proceso de integración continua sin degradar los tiempos de compilación.

\section{Conclusiones y Referencias}
\label{sec:conclusiones_referencias}

\subsection{Conclusiones}

El desarrollo e implementación de la plataforma \textbf{artesa} aporta una solución robusta y verificable al reto histórico de fragmentación operativa en micro, pequeñas y medianas empresas del sector panadero y gastronómico en El Salvador. A partir de los resultados obtenidos durante el ciclo de ingeniería de software, se derivan las siguientes conclusiones fundamentales:

\begin{enumerate}
    \item \textbf{Eficacia del Modelo Transaccional SQL:} La decisión arquitectónica de concentrar la lógica transaccional crítica en procedimientos almacenados en \textbf{PL/pgSQL} dentro de PostgreSQL 16 garantizó la consistencia absoluta de las operaciones ACID. La deducción automática de insumos a través de recetas complejas se ejecuta de manera atómica, impidiendo estados intermedios inconsistentes o corrupción de existencias ante contingencias de red.
    \item \textbf{Alineación Normativa y Fiscal Salvadoreña:} La parametrización tributaria (desglose exacto del 13\% de IVA, emisión simulada de Facturas de Consumidor Final y Comprobantes de Crédito Fiscal) y laboral (tramos de Impuesto sobre la Renta, topes y tasas de ISSS y AFP) proporciona una base técnica directamente transferible a entornos productivos reales bajo la supervisión del Ministerio de Hacienda \cite{mh_elsalvador_dte, asamblea_iva, asamblea_renta}.
    \item \textbf{Mitigación del Error de Coma Flotante (Anti-Float):} La adopción sistemática del tipo de dato \texttt{NUMERIC} en persistencia y su manipulación estricta mediante cadenas formateadas en el transporte resolvió de raíz las discrepancias centesimales clásicas en cálculos financieros acumulados.
    \item \textbf{Cultura de Calidad Continua:} La suite de 128 pruebas automatizadas ejecutadas en menos de un segundo validó el 100\% de los contratos de la API y las máquinas de estados finitos, facilitando el desarrollo colaborativo sin regresiones.
\end{enumerate}

\subsection{Trabajo Futuro}
Como líneas de extensión del proyecto se contempla la integración directa con el servicio de firma electrónica y recepción de DTE del Ministerio de Hacienda, el desarrollo de un módulo de pronóstico de demanda de materias primas mediante análisis de series temporales y la implementación de una aplicación móvil para inventarios físicos con escaneo directo de código de barras.

% Inclusión de bibliografía en formato IEEE
\bibliographystyle{IEEEtran}
\bibliography{referencias}

\section{Legalidad y Ética en el Uso de la Inteligencia Artificial}
\label{sec:legalidad_ia}

\subsection{Marco Deontológico y Normativa Académica}
En consonancia con las directrices de integridad académica de la cátedra de Programación Web (Ciclo 02-2026) y los principios deontológicos del Código de Ética de la ACM (\textit{Association for Computing Machinery}) \cite{acm_ai_ethics}, el equipo de desarrollo de \textbf{artesa} adoptó un protocolo explícito y verificable para la integración de herramientas de Inteligencia Artificial Generativa (IAG).

El uso de modelos de lenguaje de gran escala (LLMs) se define formalmente como \textbf{asistencia técnica aumentada y programación en parejas (\textit{AI Pair Programming})}, supeditada en todo momento a la supervisión, validación y autoría intelectual del equipo humano.

\subsection{Principios Rectores del Protocolo de IA}

\begin{enumerate}
    \item \textbf{Principio de Comprensión y Defensa Personal:} Cada miembro del equipo asume la responsabilidad total de las líneas de código generadas con soporte de IA. Ningún cambio es integrado al árbol principal de control de versiones (\texttt{git}) sin que el integrante responsable sea capaz de justificar su funcionamiento, trazar su flujo de ejecución y defender técnicamente su idoneidad ante el evaluador docente.
    \item \textbf{Transparencia y Trazabilidad Obligatoria:} De acuerdo con los requerimientos de la cátedra, cada interacción significativa con agentes de IA se registra formalmente en el archivo de auditoría \texttt{docs/USO\_IA.md} y en el historial de cambios del proyecto (\texttt{CHANGELOG\_AGENTE.md}), especificando fecha, operador, herramienta, alcance y verificador humano.
    \item \textbf{Confidencialidad y Sanitización de Datos:} Ningún secreto de infraestructura, clave simétrica JWT, contraseña de base de datos ni información sensible fue provista a los modelos de lenguaje. Se emplearon plantillas sanitizadas de entorno (\texttt{.env.example}) y datos de prueba sintéticos.
    \item \textbf{Verificación por Pruebas Automatizadas:} Todo código o módulo derivado de interacciones con IA fue contrastado contra la suite de pruebas unitarias y de contrato (\texttt{node --test}) antes de su aprobación final, garantizando que no se introdujeran efectos secundarios indeseados o vulnerabilidades de seguridad.
\end{enumerate}

\subsection{Registro Consolidado de Intervenciones con IA}
La Tabla~\ref{tab:registro_ia_dinamico} presenta el registro vivo de intervenciones significativas con asistentes de IA en el proyecto. Este apartado se actualiza de manera continua conforme evoluciona el desarrollo del sistema mediante el script automatizado \texttt{actualizar\_ia.js}, el cual sincroniza las bitácoras declaradas en \texttt{docs/USO\_IA.md}.

% Archivo autogenerado por actualizar_ia.js desde docs/USO_IA.md
% Fecha de actualización: 2026-11-12

\begin{table*}[htbp]
\caption{Registro Dinámico de Intervenciones con IA (Sincronizado con docs/USO\_IA.md)}
\label{tab:registro_ia_dinamico}
\centering\scriptsize
\begin{tabularx}{\textwidth}{llllXX}
\toprule
\textbf{Fecha} & \textbf{Integrante} & \textbf{Herramienta} & \textbf{Propósito / Módulo} & \textbf{Detalle de la Asistencia} & \textbf{Verificación Humana} \\
\midrule
2026-11-10 & Rodri & Codex & Separación RBAC del Kárdex con seeder idempotente, permisos backend/frontend y pruebas;... & Separación RBAC del Kárdex con seeder idempotente, permisos backend/frontend y pruebas;... & Rodri revisó y autorizó la ... \\
2026-11-10 & EmiValdeCobo & Antigravity (Gemini) & Diagnóstico y corrección de error 500 en login por columna ip\_cliente faltante & Diagnóstico y corrección de error 500 en login por columna ip\_cliente faltante: aplicac... & EmiValdeCobo \\
2026-11-10 & EmiValdeCobo & Antigravity (Gemini) & Menú POS agrupado por categorías con orden de BD, control RBAC de CCF (POS\_EMITIR\_CCF) ... & Menú POS agrupado por categorías con orden de BD, control RBAC de CCF (POS\_EMITIR\_CCF) ... & EmiValdeCobo \\
2026-11-10 & EmiValdeCobo & Antigravity (Gemini) & Reestructuración de submenús de navegación en acordeón sincronizado con ruta activa, en... & Reestructuración de submenús de navegación en acordeón sincronizado con ruta activa, en... & EmiValdeCobo \\
2026-11-10 & Rodri & Codex & Corregidas las advertencias react-hooks/exhaustive-deps en KardexPage estabilizando suc... & Corregidas las advertencias react-hooks/exhaustive-deps en KardexPage estabilizando suc... & Rodri solicitó el ajuste y ... \\
2026-11-10 & Emiliano & ChatGPT & Revisión de estructura, sistema visual, componentes compartidos y vistas SQL; rediseño ... & Revisión de estructura, sistema visual, componentes compartidos y vistas SQL; rediseño ... & frontend/src/modules/empres... \\
2026-11-10 & Emiliano & Codex & Ajuste visual responsivo del directorio de sucursales & Ajuste visual responsivo del directorio de sucursales: misma altura y estilo para búsqu... & Emiliano \\
2026-11-10 & Emiliano & Codex & Reorganización responsiva del panel de reportes con pestañas y gráfica adaptable para l... & Reorganización responsiva del panel de reportes con pestañas y gráfica adaptable para l... & Emiliano \\
2026-11-10 & EmiValdeCobo & Antigravity (Gemini) & Documentación arquitectónica y operativa en README.md & Documentación arquitectónica y operativa en README.md: generación de 5 diagramas Mermai... & EmiValdeCobo \\
2026-11-10 & EmiValdeCobo & Antigravity (Gemini) & Creación de la Documentación IEEE completa en LaTeX modular, compilador autónomo Tecton... & Creación de la Documentación IEEE completa en LaTeX modular, compilador autónomo Tecton... & EmiValdeCobo \\
2026-11-10 & EmiValdeCobo & Antigravity (Gemini) & Implementación y verificación del rediseño visual del dashboard (tarjetas redondeadas, ... & Implementación y verificación del rediseño visual del dashboard (tarjetas redondeadas, ... & EmiValdeCobo \\
2026-11-12 & Emiliano & Antigravity (Gemini) & Alineación de metadatos (Usuario, Sucursal, Hoy) con la ceja Inicio mediante items-star... & Alineación de metadatos (Usuario, Sucursal, Hoy) con la ceja Inicio mediante items-star... & Emiliano \\
\bottomrule
\end{tabularx}
\end{table*}


\subsection{Conformidad sobre Propiedad Intelectual y Licenciamiento}
El código resultante se encuentra bajo licencia académica y de código abierto compatible con el ecosistema de Node.js, React y PostgreSQL. Las sugerencias de código provistas por los modelos de IA fueron asimiladas como estructuras sintácticas estándar de la industria, libres de reclamaciones de derechos de autor de terceros y totalmente auditadas por los estudiantes autores de este documento.


\end{document}
